\chapter{多进程编程}
\label{cha:多进程编程}
\section{进程}
\label{sec:进程}
Python3中用到\codeinline{multiprocessing}模块来实现多进程编程，进程是操作系统能够进行运算调度的最小单位，是程序的一次执行过程，是一个动态的过程：有它自己的地址空间，有它自己的内存，有它自己的数据，有它自己的打开的文件等等。进程的切换会消耗大量的资源，所以进程间的通信相对复杂。

\begin{lstlisting}[style=demostyle,numbers=none]
Process([group [, target [, name [, args [, kwargs]]]]])
'''
group: 线程组，目前还没有实现，库引用中应该写None
target: 要执行的方法
name: 进程名
args/kwargs: 方法的参数
'''
start() # 启动进程
join([timeout]) # 等待进程终止
\end{lstlisting}

\begin{lstlisting}[style=pythonstyle,numbers=left]
# 不推荐的写法
import multiprocessing

def func(arg1, arg2):
    for i in range(5):
        print('multiprocessing: %s, %s %s' % (i, arg1+i, i+arg2))

p = multiprocessing.Process(target=func, args=(1,2))
p.start()
\end{lstlisting}

一个好的方式是继承\codeinline{multiprocessing.Process}类，然后重写\codeinline{run}方法，这样可以更加灵活、直观。

\begin{lstlisting}[style=pythonstyle]
# 推荐的写法
## 上面的方法虽然用起来简单，但是不推荐使用，因为这样写不利于维护，韩式推荐使用继承的方式来创建进程
import multiprocessing

def func(arg1, arg2):
    for i in range(5):
        print('multiprocessing: %s, %s %s' % (i, arg1+i, i+arg2))

class Myprocess(multiprocessing.Process):
    def __init__(self, func, args, pname=''):
        super(Myprocess, self).__init__()
        self.func = func    # func 存放的是要执行的函数
        self.args = args    # args 存放的函数的参数
        self.pname = pname  # pname 存放的进程的名字，进程的名字是可选的，
                            # 如果不传入进程名字，进程会自动分配一个名字

    ## 执行进程的具体逻辑
    ## 通过重写run方法的方式可以自定义进程具体的执行逻辑，
    ## 相比上面的方法， 这种方法更加灵活、直观。
    def run(self):
        print("Runing Process: ",self.pname)
        print('run')
        self.func(*self.args)
        print(current_process().name, "end")

if __name__ == '__main__':
    p = Myprocess(func, (1, 2), 'myprocess')
    p.start()
    p.join()
\end{lstlisting}

\begin{lstlisting}[style=demostyle,numbers=none]
$ python  Myprocess.py
Runing Process:  myprocess
Myprocess-1 start
run
multiprocessing: 0, 1 2
multiprocessing: 1, 2 3
multiprocessing: 2, 3 4
multiprocessing: 3, 4 5
multiprocessing: 4, 5 6
Myprocess-1 stop
\end{lstlisting}

\section{进程间通信}
\label{sec:进程间通信}
\begin{Tips}[进程与线程的区别]
    与线程不同，线程如同车间里的工人，工人之间可以直接使用车间里的任意资源；

    但进程表示车间本身，车间与车间是相互独立的，各自有自身的空间让自己的线程完成具体的子任务，车间之间的信息与资源是不互通的，如何在不同进程中传递信息是进程在使用过程中需要解决的问题。
\end{Tips}

\subsection{Queue}
\label{subsec:进程通信Queue}

\begin{lstlisting}[style=pythonstyle,numbers=left]
# 利用Queue实现生产者消费者模型
# 利用Queue实现生产者消费者模型
import os
import time
from multiprocessing import Process, Queue

def producer(q):
    for i in range(5):
        print(f'\033[32m生产者: {os.getpid()} 生产了商品{i}\033[0m')
        # 将商品放入队列
        q.put(i)
        time.sleep(1)

def consumer(q):
    while True:
        try:
            # 从队列中获取商品
            i = q.get(timeout=5)
            print(f'\033[95m消费者: {os.getpid()} 消费了商品{i}\033[0m')
        except:
            break

if __name__ == '__main__':
    q = Queue()
    p = Process(target=producer, args=(q,))
    c = Process(target=consumer, args=(q,))
    p.start() # 生产者进程开始生产商品
    c.start() # 消费者进程开始消费商品
    p.join()  # 主进程挂起，等待子进程producer
    c.join()  # 主进程挂起，等待子进程consumer
    print('主进程结束')
\end{lstlisting}

\begin{lstlisting}[style=demostyle,numbers=none]
$ python queue.py
生产者: 19896 生产了商品0
消费者: 5856 消费了商品0
生产者: 19896 生产了商品1
消费者: 5856 消费了商品1
生产者: 19896 生产了商品2
消费者: 5856 消费了商品2
生产者: 19896 生产了商品3
消费者: 5856 消费了商品3
生产者: 19896 生产了商品4
消费者: 5856 消费了商品4
主进程结束
\end{lstlisting}


\section{进程池}
\label{sec:进程池}
进程池是一种能够提供多个进程的服务，进程池中的进程可以执行任务，当任务执行完毕后，进程会自动回到进程池中，等待下一个任务。

\begin{lstlisting}[style=demostyle,numbers=none]
Pool([processes [, initializer [, initargs [, maxtasksperchild [, context]]]]])
'''
processes: 进程池中的进程数
initializer: 初始化函数
initargs: 初始化函数的参数
maxtasksperchild: 每个子进程最多执行的任务数
context: 上下文
'''
apply(func[, args[, kwds]]) # 同步执行
apply_async(func[, args[, kwds]]) # 异步执行
map(func, iterable[, chunksize]) # 同步执行
map_async(func, iterable[, chunksize]) # 异步执行
\end{lstlisting}

\subsection{同步执行}
\label{subsec:同步执行}
\begin{lstlisting}[style=pythonstyle,numbers=left]
# 进程同步
# 演示如何使用 multiprocessing.Pool 来创建进程池，
# 以及如何使用 apply 和 map 方法来执行函数

from multiprocessing import Pool
import os
import time

def square(x):
    print(f'Process {os.getpid()} calculating square of {x}')
    time.sleep(1)  # 模拟耗时操作
    return x * x

def worker(num):
    print(f'Process {os.getpid()} doubling {num}')
    time.sleep(1)  # 模拟耗时操作
    return num * 2

if __name__ == "__main__":
    with Pool(processes=4) as pool: # 创建进程池
        # 使用 apply 同步执行
        print("Starting apply()")
        ## apply方法会阻塞主进程，直到子进程执行完毕并返回结果
        result_apply = pool.apply(square, (10,))
        print(f'Result from apply: {result_apply}')

        # 使用 map 同步执行
        print("Starting map()")
        numbers = [1, 2, 3, 4, 5]
        result_map = pool.map(worker, numbers)
        print(f'Results from map: {result_map}')
\end{lstlisting}

\begin{lstlisting}[style=demostyle,numbers=none]
$ python pool.py
Starting apply()
Process 17668 calculating square of 10
Result from apply: 100
Starting map()
Process 4084 doubling 1
Process 32672 doubling 2
Process 12144 doubling 3
Process 17668 doubling 4
Process 12144 doubling 5
Results from map: [2, 4, 6, 8, 10]
\end{lstlisting}

\subsection{异步执行}
\label{subsec:异步执行}

\begin{lstlisting}[style=pythonstyle,numbers=left]
# 进程异步
# 演示如何使用 multiprocessing.Pool 来创建进程池，
# 以及如何使用 apply_async 和 map_async 方法来异步执行函数

from multiprocessing import Pool
import os
import time

def square(x):
    print(f'Process {os.getpid()} calculating square of {x}')
    time.sleep(1)  # 模拟耗时操作
    return x * x

def worker(num):
    print(f'Process {os.getpid()} doubling {num}')
    time.sleep(1)  # 模拟耗时操作
    return num * 2

if __name__ == "__main__":
    with Pool(processes=4) as pool:
        # 使用 apply_async 异步执行
        print("Starting apply_async()")
        result_apply_async = pool.apply_async(square, (10,))
        
        # 在等待结果的同时，可以做其他事情
        print("Doing some other work while waiting for apply_async result")
        time.sleep(2)  # 模拟其他工作耗时
        print(f'Result from apply_async: {result_apply_async.get()}')  # 获取结果

        # 使用 map_async 异步执行
        print("Starting map_async()")
        numbers = [1, 2, 3, 4, 5]
        result_map_async = pool.map_async(worker, numbers)
        
        # 在等待结果的同时，可以做其他事情
        print("Doing some other work while waiting for map_async results")
        time.sleep(2)  # 模拟其他工作耗时
        print(f'Results from map_async: {result_map_async.get()}')  # 获取结果
\end{lstlisting}

\begin{lstlisting}[style=demostyle,numbers=none]
$ python pool_async.py
Starting apply_async()
Doing some other work while waiting for apply_async result
Process 35528 calculating square of 10
Result from apply_async: 100
Starting map_async()
Doing some other work while waiting for map_async results
Process 6792 doubling 1
Process 36048 doubling 2
Process 36592 doubling 3
Process 35528 doubling 4
Process 35528 doubling 5
Results from map_async: [2, 4, 6, 8, 10]
\end{lstlisting}






